% Source PDF: source/普通高考/2011/2011安徽文.pdf, page 1, question 12.
% PDF embedded bitmap attempt: pdfimages -list returned no image objects; the flowchart is vector artwork.
% Grid evidence: tmp/review_flowchart_2011_anhui_liberal/grid/page-1-grid100.png
% (100px coarse + 50px fine) and q12-block-grid50.png (50px coarse + 25px fine).
% Original-side comparison crop: tmp/review_flowchart_2011_anhui_liberal/crop/q12_original_final.png,
% taken from the ungridded page render with a safety border.
% Elements: rounded start/end terminals, process boxes, two decision/output branches,
% one T>105? decision diamond, output parallelogram, and directed loop connectors.
% Node-edge list: start->(T=0,k=0)->(T=T+k)->(T>105?); yes->输出k->结束;
% no->(k=k+1)->loop-back to the stem before T=T+k. No other edges or dashed segments.
% The loop arrow enters the main stem from the right with one horizontal arrowhead;
% the output path remains a separate downward branch.
\begin{tikzpicture}[
  >=Stealth,
  x=1cm,y=1cm,
  line width=0.45pt,
  line cap=round,line join=round,
  font=\normalsize,
  flowtext/.style={anchor=center,align=center,inner sep=0pt,
    execute at begin node={\setlength{\parindent}{0pt}\noindent}},
  flowbox/.style={flowtext,draw,minimum height=.68cm,inner xsep=4pt,inner ysep=2pt},
  terminal/.style={flowbox,rounded corners=.16cm,minimum width=1.25cm},
  process/.style={flowbox},
  io/.style={flowbox,shape=trapezium,trapezium left angle=72,trapezium right angle=108,
    minimum height=.74cm},
  decision/.style={flowbox,diamond,aspect=2.8,minimum width=3.10cm,minimum height=.90cm,
    inner sep=0pt}
]
  % Keep a small right margin so the k=k+1 frame is not clipped.
  \path[use as bounding box] (-3.55,-6.25) rectangle (4.85,0.42);

  \node[terminal] (start) at (0,0) {开始};
  \node[process,minimum width=3.10cm] (init) at (0,-1.00) {$T=0,k=0$};
  \node[process,minimum width=2.40cm] (sum) at (0,-2.30) {$T=T+k$};
  \node[decision] (test) at (0,-3.45) {$T>105?$};
  \node[process,minimum width=2.20cm] (inc) at (3.50,-3.45) {$k=k+1$};
  \node[io,minimum width=1.80cm] (output) at (0,-4.58) {输出 $k$};
  \node[terminal] (finish) at (0,-5.70) {结束};
  % Leave at least 2.5 mm around the shared stem and the neighboring boxes.
  \coordinate (merge-level) at (0,-1.65);
  \coordinate (return-entry) at (0.35,-1.65);

  % Main chain and terminating true branch.
  \draw[->] (start.south) -- (init.north);
  \draw (init.south) -- (merge-level);
  \draw[->] (merge-level) -- (sum.north);
  \draw[->] (sum.south) -- (test.north);
  \draw[->] (test.south) -- node[right=2pt] {是} (output.north);
  \draw[->] (output.south) -- (finish.north);

  % False branch updates k, then returns up to the stem before T=T+k.
  \draw[->] (test.east) -- node[above=2pt,pos=.38] {否} (inc.west);
  \coordinate (loop-right) at (3.50,-1.78);
  % The return leg leaves the top of the update box, avoiding a line through
  % the node interior before turning left into the shared stem.
  \draw (inc.north) -- (loop-right);
  % The directed return path ends at the shared stem entry.
  \draw[->] (loop-right) -- (return-entry) -- (merge-level);
  
\end{tikzpicture}
